\section{Missing transverse momentum and object definition}
\label{sec:ch6_met}

\subsection{Reconstruction}

The incoming partons at the LHC carry little transverse momentum, so the vector sum of the transverse momenta of all the particles is expected to be approximately zero.
An imbalance in the visible transverse momentum can therefore indicate undetected particles, such as neutrinos or weakly interacting BSM particles.
In this analysis, the neutrino produced in the LQ process and the tau-decay neutrino can together give rise to sizeable missing transverse momentum ($\mpt$).
$\mpt$ is reconstructed as the negative vector sum of the calibrated visible contributions, with its magnitude denoted by $\met$~\cite{ATLAS2025METPerformance}:
\begin{align}
    \mpt
    &=-\left(
    \sum_e\vec p_{\mathrm{T}}^{\,e}
    +\sum_\mu\vec p_{\mathrm{T}}^{\,\mu}
    +\sum_{\tau_{\mathrm{had}}}\vec p_{\mathrm{T}}^{\,\tau}
    +\sum_j\vec p_{\mathrm{T}}^{\,j}
    +\vec p_{\mathrm{T}}^{\,\mathrm{soft}}
    \right),
    \label{eq:ch6_met_vector}\\
    \met&=\left|\mpt\right|.
    \label{eq:ch6_met_magnitude}
\end{align}
Here, baseline objects passing the OLR procedure is used. And two types of contribution are included:
\begin{itemize}
    \item \textbf{hard term}: The electron, muon, hadronic-tau and jet contributions selected for $\mpt$ reconstruction.
    The jet contribution is selected using the \texttt{Tight} $\mpt$ working point to suppress pile-up jets.
    This requires jet $\pT>20~\GeV$ for $|\eta|<2.5$ and $\pT>30~\GeV$ for $2.5\leq|\eta|<4.5$.
    \item \textbf{track soft term}: The vector sum of transverse momenta of selected PV tracks not used in the hard term.
\end{itemize}
